\chapter{书写规范与打印要求}
\section{文字和字数}
一般用汉语简化文字书写，字数在1.2万字左右，报告（内容）或软件说明书，字数在1万字左右。

\section{书写}
论文一律由本人在计算机上输入、编排并打印在A4幅面白纸上。毕业论文（设计）前置部分（即正文之前）一律用单面印刷，正文部分开始双面印刷。致谢和附录部分应单面起页双面印刷（如正文结束页为单页，则单数页背面不加页眉和页码，致谢单面起页，如致谢为单页，其背面亦不加页眉和页码）。

\section{字体和字号}
论~~文~~题目：\hspace{2cm}小2号黑体

章~~~~标~~~~题：\hspace{2cm}3号黑体

节~~~~标~~~~题：\hspace{2cm}小4号黑体

条~~~~标~~~~题：\hspace{2cm}小4号黑体

正~~~~~~~~~~~~文：\hspace{2cm}小4号宋体

页~~~~~~~~~~~~码：\hspace{2cm}5号宋体

数字和字母：\hspace{2cm}Times New Roman体

\section{封面}
毕业论文（设计）封面规范（见样张1），论文封皮一律采用白色铜版纸，封皮大小为A4规格。

\section{原创性声明及授权书}
见样张2。

\section{论文页面设置}
\subsection{页眉}
页眉为\uline{~~~~~~~~湖南大学本科生毕业论文（设计）~~~~~~~~}。

\subsection{页边距}
上边距：30mm；下边距：25mm；左边距：30mm；右边距：20mm；行间距为1.5倍行距。

\subsection{页码的书写要求}
论文页码从绪论部分开始，至附录，用阿拉伯数字连续编排，页码位于页脚居中。封面、摘要和目录不编入论文页码；摘要和目录用罗马数字单独编页码。

\section{摘要}
\subsection{中文摘要}
中文摘要包括：论文题目（小3号黑体）、“摘要”字样（3号黑体）、摘要正文(小4号宋体)和关键词。

摘要正文后下空一行打印“关键词”三字（4号黑体），关键词一般为3～5个，每一关键词之间用分号分开，最后一个关键词后不打标点符号。

\subsection{英文摘要}
英文摘要另起一页，其内容及关键词应与中文摘要一致，并要符合英语语法，语句通顺，文字流畅。

英文题目的字样为：The title（小3号 Times New Roman 加粗）

英文摘要字样为：Abstract(3号 Times New Roman 加粗)

然后隔行书写摘要的文字部分。（字体为小4号Times New Roman）

摘要正文后下空一行打印关键词（4号Times New Roman加粗）：key word1；key word2；（关键词3－5个，小4号Times New Roman加粗）

英文和汉语拼音一律为Times New Roman体，字号与中文摘要相同。

\section{目录}
专业目录的三级标题，建议按（1……、1.1……、1.1.1……）的格式编写，目录中各章题序的阿拉伯数字用Times New Roman体，第一级标题用小4号黑体，其余用小4号宋体。目录的打印范例。

\section{论文正文}
\subsection{章节及各章标题}
论文正文分章节撰写，每章应另起一页。各章标题要突出重点、简明扼要。字数一般在15字以内，不得使用标点符号。标题中尽量不采用英文缩写词，对必须采用者，应使用本行业的通用缩写词。

\subsection{层次}
层次以少为宜，根据实际需要选择。正文层次的编排和代号要求统一，层次为章（如“1”）（3号黑体）、节（如“1.1”）、条（如“1.1.1”）、款（如“（1）”）。层次用到哪一层次视需要而定，若节后无需“条”时可直接列“款”、“项”。“节”、“条”的段前、段后各设为0.5行。

\section{引用文献}
引用文献标示方式应全文统一，并采用所在学科领域内通用的方式，所引文献编号用阿拉伯数字置于方括号中，用上标的形式置于所引内容最末句的右上角，如：“…成果$^{[1]}$”，引用文献应与文中标注一致。几处地方引用同一个文献时，文中标注按第一次出现的序号。当提及的参考文献为文中直接说明时，其序号应该用阿拉伯数字与正文排齐，如“由文献[8, 10-14]可知”。

不得将引用文献标示置于各级标题处。

\section{名词术语}
科技名词术语及设备、元件的名称，应采用国家标准或部颁标准中规定的术语或名称。标准中未规定的术语要采用行业通用术语或名称。全文名词术语必须统一。一些特殊名词或新名词应在适当位置加以说明或注解。

采用英语缩写词时，除本行业广泛应用的通用缩写词外，文中第一次出现的缩写词应该用括号注明英文全文。

\section{物理量名称、符号与计量单位}
\subsection{物理量的名称和符号}
物理量的名称和符号应符合GB3100~3102-86的规定。论文中某一量的名称和符号应统一。

\subsection{物理量计量单位}
物理量计量单位及符号应按国务院1984年发布的《中华人民共和国法定计量单位》及GB3100\~3102执行，不得使用非法定计量单位及符号。计量单位符号，除用人名命名的单位第一个字母用大写之外，其它一律用小写字母。

非物理量单位（如件、台、人、元、次等）可以采用汉字与单位符号混写的方式，如“万t·km”。

文稿叙述中不定数字之后允许用中文计量单位符号，如“几千克至1000kg”。

表达时刻时应采用中文计量单位，如“上午8点3刻”，不能写成“8h45min”。

计量单位符号一律用正体。

\section{外文字母的正、斜体用法}
物理量符号、物理常量、变量符号用斜体，计量单位等符号均用正体。

\section{数字}
除习惯用中文数字表示的以外，一般均采用阿拉伯数字。年份一概写全数，如2003年不能写成03年。

\section{公式}
公式应另起一行写在稿纸中央，公式和编号之间不加虚线。公式较长时最好在等号“＝”处转行，如难实现，则可在＋、－、×、÷运算符号处转行，运算符号应写在转行后的行首，公式的编号用圆括号括起来放在公式右边行末。

公式序号按章编排，如第一章第一个公式序号为“(1.1)”，附录A中的第一个公式为“(A1)”等。

文中引用公式时，一般用“见式(1.1)”或“由公式(1.1)”。

公式中用斜线表示“除”的关系时应采用括号，以免含糊不清，如a/(bcosx)。通常“乘”的关系在前，如acosx/b而不写成(a/b)cosx。

\section{表格}
每个表格应有的表序和表题。并应在文中进行说明，例如：“如表1.1”。

表序一般按章编排，如第一章第一个插表的序号为“表1.1”等。表序与表名之间空一格，表名中不允许使用标点符号，表名后不加标点。表序与表名置于表上居中（5号黑体加粗，数字和字母为5号Times New Roman体加粗）。

表头设计应简单明了，尽量不用斜线。表头与表格为一整体，不得拆开排写于两页。

全表如用同一单位，将单位符号移至表头右上角。

表中数据应正确无误，书写清楚。数字空缺的格内加“－”字线（占2个数字），不允许用“²”、“同上”之类的写法。

表内文字说明（5号宋体），起行空一格、转行顶格、句末不加标点。

表中若有附注时，用小5号宋体，写在表的下方，句末加标点。仅有一条附注时写成：注：；

有多条附注时，附注各项的序号一律用阿拉伯数字，例如：注1。

\section{图}
毕业设计的插图应与文字紧密配合，文图相符，技术内容正确。选图要力求精练。

\subsection{制图标准}
插图应符合国家标准及专业标准。

机械工程图：采用第一角投影法，严格按照GB4457~4460-84，GB131-83《机械制图》标准规定。

电气图：图形符号、文字符号等应符合有关标准的规定。

流程图：原则上应采用结构化程序并正确运用流程框图。

对无规定符号的图形应采用该行业的常用画法。

\subsection{图题及图中说明}
每幅插图均应有图题（由图号和图名组成）。图号按章编排，如第一章第一图的图号为“图1.1”等。图题置于图下，用5号宋体加粗。有图注或其他说明时应置于图题之上，用小5号宋体。图名在图号之后空一格排写。引用图应说明出处，在图题右上角加引用文献号。图中若有分图时，分图号用(a)、(b)等置于分图之下。

图中各部分说明应采用中文（引用的外文图除外）或数字符号，各项文字说明置于图题之上（有分图题者，置于分图题之上）。
\subsection{插图编排}
插图与其图题为一个整体，不得拆开排写于两页。插图处的该页空白不够排写该图整体时，可将其后文字部分提前排写，将图移至次页最前面。
\subsection{坐标与坐标单位}
对坐标轴必须进行说明，有数字标注的坐标图，必须注明坐标单位。
\subsection{论文原件中照片图及插图}
毕业设计(论文)原件中的照片图应是直接用数码相机拍照的照片，或是原版照片粘贴，不得采用复印方式。照片可为黑白或彩色，应主题突出、层次分明、清晰整洁、反差适中。照片采用光面相纸，不宜用布纹相纸。对金相显微组织照片必须注明放大倍数。

\section{注释}
毕业设计（论文）中有个别名词或情况需要解释时，可加注说明，注释可用页末注（将注文放在加注页稿纸的下端）或篇末注（将全部注文集中在文章末尾），而不用行中注（夹在正文中的注）。若在同一页中有两个以上的注时，按各注出现的先后编列注号，注释只限于写在注释符号出现的同页，不得隔页。

\section{参考文献}
参考文献的著录均应符合国家有关标准（按GB7714—2015《信息与文献参考文献著录规则》执行）。详见附件《参考文献示例》。

关于参考文献的未尽事项可参见国家标准《信息与文献参考文献著录规则》（GB7714－2015）。

\section{……}
其他部分请自行阅读学院颁布的撰写规范文件。